% Part: sets-functions-relations
% Chapter: relations
% Section: equivalence-relations

\documentclass[../../../include/open-logic-section]{subfiles}

\begin{document}

\olfileid{sfr}{rel}{eqv}

\olsection{Relasi Ekuivalensi}

Relasi identitas ing sawijining himpunan iku refleksif, simetris,
lan transitif. Relasi~$R$ kang duwe telung sipat iki kerep
banget ditemoni.

\begin{defn}[Relasi ekuivalensi]
Relasi $R \subseteq A^2$ kang refleksif, simetris,
lan transitif diarani \emph{relasi ekuivalensi}.
!!^{element}s $x$ lan $y$ saka~$A$ diarani
\emph{ekuivalen tumrap $R$} yen~$Rxy$.
\end{defn}

Relasi ekuivalensi nuwuhake gagasan \emph{kelas ekuivalensi}.
Relasi ekuivalensi "merang" domain dadi bageyan-bageyan
partisi kang beda. Ing saben bageyan, kabeh objek gegayutan
siji lan sijine; lan ora ana objek saka bageyan beda kang
gegayutan siji lan sijine. Kadhangkala migunani yen
ngrembug bageyan-bageyan iki kanthi \emph{langsung}.
Kanggo iku, kita ngenalake definisi iki:

\begin{defn}\ollabel{def:equivalenceclass}
Umpamakna $R \subseteq A^2$ iku relasi ekuivalensi.
Kanggo saben $x \in A$, \emph{kelas ekuivalensi} saka $x$
ing~$A$ yaiku himpunan
$\equivrep{x}{R}= \Setabs{y \in A}{Rxy}$.
\emph{Kuosien} saka $A$ tumrap~$R$ yaiku
$\equivclass{A}{R} = \Setabs{\equivrep{x}{R}}{x \in A}$,
yaiku himpunan kelas-kelas ekuivalensi iki.
\end{defn}

Asil sabanjure mbenerake definisi kelas ekuivalensi kanthi
mbuktekake yen kelas-kelas ekuivalensi pancen bageyan-bageyan
partisi saka~$A$:

\begin{prop}
Yen $R \subseteq A^2$ iku relasi ekuivalensi,
mula $Rxy$ yen lan mung yen
$\equivrep{x}{R} = \equivrep{y}{R}$.
\end{prop}

\begin{proof}
Kanggo arah kiwa menyang tengen, umpamakna $Rxy$,
lan tetepna $z \in\equivrep{x}{R}$.
Miturut definisi, banjur $Rxz$.
Amarga $R$ iku relasi ekuivalensi, mula $Ryz$.
(Kanthi luwih rinci: amarga $Rxy$ lan~$R$ simetris,
kita oleh $Ryx$, lan amarga $Rxz$ lan~$R$ transitif,
kita oleh~$Ryz$.) Mula $z \in \equivrep{y}{R}$.
Kanthi umum, $\equivrep{x}{R} \subseteq \equivrep{y}{R}$.
Nanging kanthi cara kang persis padha,
$\equivrep{y}{R} \subseteq \equivrep{x}{R}$.
Mula $\equivrep{x}{R} =\equivrep{y}{R}$,
miturut ekstensionalitas.

Kanggo arah tengen menyang kiwa, umpamakna
$\equivrep{x}{R} =\equivrep{y}{R}$.
Amarga $R$ refleksif, $Ryy$, mula
$y \in\equivrep{y}{R}$. Dadi uga
$y \in \equivrep{x}{R}$, amarga anggepan yen
$\equivrep{x}{R} = \equivrep{y}{R}$.
Mula $Rxy$.
\end{proof}

\begin{ex}
Tuladha relasi ekuivalensi kang becik asale saka aritmetika modular.
Kanggo saben $a$, $b$, lan $n \in \PosInt$,
kita nyebut $a \equiv_n b$ yen lan mung yen
mbagi $a$ nganggo~$n$ menehi sisa kang padha
karo mbagi $b$ nganggo~$n$.
(Kanthi lambang: $a \equiv_n b$ yen lan mung yen
ana $k \in\Int$ kang ndadekake $a - b = kn$.)
Saiki, $\equiv_n$ iku relasi ekuivalensi kanggo saben~$n$.
Lan ana persis $n$ kelas ekuivalensi beda kang kawangun
dening~$\equiv_n$; yaiku, $\equivclass{\Nat}{\equiv_n}$
duwe $n$ !!{element}s.
Kelas-kelas iki yaiku: himpunan wilangan kang bisa dibagi $n$
tanpa sisa, yaiku $\equivrep{0}{\equiv_n}$;
himpunan wilangan kang dibagi $n$ menehi sisa~$1$,
yaiku $\equivrep{1}{\equiv_n}$; \ldots;
lan himpunan wilangan kang dibagi~$n$ menehi sisa~$n-1$,
yaiku~$\equivrep{n-1}{\equiv_n}$.
\end{ex}

\begin{prob}
Tuduhna yen $\equiv_n$ iku relasi ekuivalensi kanggo saben
$n \in\PosInt$, lan yen $\equivclass{\Nat}{\equiv_n}$
duwe persis $n$ anggota.
\end{prob}

\end{document}
